\section{Implementation and technology choices}

\subsection{Technology stack}
\begin{table}[htbp]
\centering
\small
\begin{tabularx}{\textwidth}{@{}p{31mm}p{35mm}X@{}}
\toprule
\textbf{Area} & \textbf{Technology} & \textbf{Reason} \\
\midrule
Application & Flutter and Dart & Typed cross-platform UI, mature Android
tooling, immutable value models, and testable widgets \\
Local database & SQLite through \code{sqflite} & Transactional, versioned,
offline application records on Android \\
Location & \code{geolocator} & Permission, service-state, position stream, and
foreground-notification settings behind a contract \\
Camera & \code{image_picker} & Native camera handoff with a small adapter surface \\
Filesystem & \code{path_provider}, \code{path} & App-controlled support and cache
locations without paths in widgets \\
Export and share & Flutter image codecs, \code{share_plus} & Deterministic PNG
generation and the Android system share sheet \\
Optional remote & HTTP plus Supabase Edge Functions & Small REST boundary,
anonymous expiring records, and no SDK in the domain \\
Testing & \code{flutter_test} and deterministic fakes & Logic, widget states,
navigation, serialization, and failure behavior without devices or a server \\
Report & XeLaTeX and \code{latexmk} & Reproducible A4 output, modular source,
stable tables, diagrams, references, and automated reruns \\
\bottomrule
\end{tabularx}
\caption{Principal technologies and their roles.}
\end{table}



Flutter provides one strongly typed UI and domain codebase while retaining
Android device integration. Material 3 supplies accessible interaction
primitives, but the final surfaces use a custom solid-color identity. No large
state-management framework was introduced; feature controllers and immutable
snapshots keep the dependency surface small and inspectable.

\subsection{Local persistence and recovery}
SQLite is opened through a small database class with an explicit schema version.
Repositories serialize active sessions, summaries, story documents, and
achievements. Destination content stays in versioned JSON because it is bundled,
reviewable, and deterministic. Photos are files because large binary payloads do
not belong in normal relational rows.

The active run is checkpointed incrementally after meaningful state changes.
When the application restarts, the main navigation checks for a checkpoint and
offers Resume or Discard. Repository operations expose failures to presentation
as recoverable states rather than allowing uncaught initialization or storage
exceptions to leave a blank screen.

\subsection{Consistency during deletion}
Deleting a run is a cascade across story documents, the run summary, and retained
photos. The orchestration removes related editable stories first, attempts the
run deletion, restores stories if that record operation fails, and performs file
cleanup only after repository state is consistent. Photo-cleanup failure is
reported separately because the user record has already been deleted.

\subsection{Optional online sharing}
Anonymous link sharing is implemented behind \code{LiveShareService}. The client
sends a deliberately limited payload to a Supabase Edge Function; the function
validates the request, stores an expiring record, and returns a share URL. The
anonymous application key may be compiled into the client, but the service-role
key remains only in the function environment. Core operation never waits for
this server. Before uploading, the UI names the shared coordinates, run time,
captions, and evidence photos and requires explicit confirmation.

There are three valid run modes:
\begin{itemize}
  \item \textbf{Offline production mode:} no server and no remote definitions;
  \item \textbf{Hosted sharing mode:} a deployed Supabase project and function;
  \item \textbf{Local sharing development:} Supabase CLI plus a Docker-compatible
        runtime, with the Android emulator addressing the host at
        \code{10.0.2.2} rather than \code{localhost}.
\end{itemize}

\subsection{Platform scope}
The release target is Android API 24 or later. The production entry point opens a
mobile SQLite database and therefore must be run on an Android device or
emulator. A browser entry point exists only as a focused Story Studio development
harness. This distinction explains the reported \code{databaseFactory not
initialized} screen on Windows/Edge: the wrong platform was selected before any
normal user-flow screen could load. It is not evidence that a backend is missing.
